\documentclass{article}


%%%%%%%%%%%%%%%%%%%%%%%%%%%%%%%%%%%%%%%%%
%%%%%%%%%%% Premiers packages %%%%%%%%%%%
%%%%%%%%%%%%%%%%%%%%%%%%%%%%%%%%%%%%%%%%%

% Package de gestion des caractères.
\usepackage[T1]{fontenc}

% Package pour avoir les bons espacements, dépendamment de la langue utilisée.

\usepackage[french]{babel}
%\usepackage[british]{babel}

%%%%%%%%%%%%%%%%%%%%%%%%%%%%%%%%%%%%%%%%%%%%%%%%%%%%%%%
%%%%%%%%%%% Modification de la mise en page %%%%%%%%%%%
%%%%%%%%%%%%%%%%%%%%%%%%%%%%%%%%%%%%%%%%%%%%%%%%%%%%%%%

% Package de gestion des interlignes.
\usepackage{setspace}

% Package de gestion des figures. 
\usepackage{subfigure}

% Package de gestion des images et des légendes.
\usepackage{graphicx}
\usepackage[export]{adjustbox}
\usepackage[font=small,labelfont=bf]{caption}

% Package de gestion des tableaux.
\usepackage{tabularx}
\usepackage{colortbl}

% Package de gestion des listes.
\usepackage{enumitem}

% Package de gestion du titre.
\usepackage{titling}
\usepackage{hyperref}
% code
\usepackage{listings}

\usepackage[
  paperwidth=21cm,
  paperheight=29.7cm,
  left=2cm,
  right=2cm,
  top=2cm,
  bottom=2cm,
  headheight=14pt,
  headsep=0.5cm,
  includehead,
  footskip=1cm
]{geometry}
% % Mise en page générale :
% \setlength{\paperwidth}{21cm}		% Largeur du papier
% \setlength{\paperheight}{29.7cm}	% Hauteur du papier

% % En haut :
% \setlength{\hoffset}{-0.54cm}		% Marge d'en haut = 1 in + \hoffset
% \setlength{\topmargin}{0cm}			% Marge de haut de page
% \setlength{\headheight}{0cm} 		% Hauteur du haut de page
% \setlength{\headsep}{0cm} 			% Entre le haut de page et le texte

% % En bas : 
% \setlength{\footskip}{1cm} 			% Bas de page + séparation

% % À droite, à gauche : 
% \setlength{\voffset}{-0.54cm} 		% Marge de gauche = 1 in + \voffset
% \setlength{\oddsidemargin}{0cm} 	% Marge gauche EN PLUS sur pages impaires
% \setlength{\evensidemargin}{0cm} 	% Marge droite EN PLUS sur pages paires

% % Zone de texte :
% \setlength{\textheight}{25.7cm} 	% Hauteur de la zone de texte (25cm)
% \setlength{\textwidth}{17cm} 		% Largeur de la zone de texte

% Permet de commencer la numérotation à 1.
\renewcommand{\thesection}{\arabic{section}}

% Interlignes à 1,5 plutôt que 1.
\renewcommand{\arraystretch}{1.5}
\onehalfspacing

%%%%%%%%%%%%%%%%%%%%%%%%%%%%%%%%%%%%%
%%%%%%%%%%% Mathématiques %%%%%%%%%%%
%%%%%%%%%%%%%%%%%%%%%%%%%%%%%%%%%%%%%

% Package général de gestion des maths.
\usepackage{amsmath}
% Packages de gestion des symboles des maths.
\usepackage{amssymb}
\usepackage{stmaryrd}
\usepackage{mathabx}
% Package de gestion des théorèmes.
\usepackage{amsthm}

\newtheorem{theorem}{Théorème}
\newtheorem{lemma}[theorem]{Lemme}
\newtheorem{conjecture}[theorem]{Conjecture}
\newtheorem{definition}[theorem]{Définition}
\newtheorem{proposition}[theorem]{Proposition}
\newtheorem{remark}[theorem]{Remarque}
\newtheorem{example}[theorem]{Exemple}

% Pour l'Expo d'ES
\newtheorem{image}{Image}[subsection]



\newtheoremstyle{exo}
  {\topsep} % space above
  {\topsep}   % space below
  {}                   % body font: empty = upright
  {}                   % indent
  {\bfseries}
  {}
  {.5em}
  {\thmname{#1}\thmnumber{ #2}\thmnote{ --- \normalfont\itshape #3}}   % heading font, punctuation, space after heading, heading spec
% 

% \newtheoremstyle{exo}
%   {\topsep}{\topsep}
%   {\normalfont}                 % upright body
%   {0pt}
%   {\bfseries}                   % heading font ("Exercice 2")
%   {}                            % no punctuation after the heading
%   {}%\newline}                    % body starts on the next line
%   % {}
%   {\thmname{#1}\thmnumber{ #2}\thmnote{ --- \normalfont\itshape #3}}
\theoremstyle{exo}
\newtheorem{exercise}{Exercice}[section]
\newtheorem*{solution}{Solution}

\theoremstyle{plain}
% Package de gestion des figures géométriques en maths.
\usepackage{tikz} 

%New colors defined below
\definecolor{codegreen}{rgb}{0,0.6,0}
\definecolor{codegray}{rgb}{0.5,0.5,0.5}
\definecolor{codepurple}{rgb}{0.58,0,0.82}
\definecolor{backcolour}{rgb}{0.95,0.95,0.92}

%Code listing style named "mystyle"
\lstdefinestyle{mystyle}{
  backgroundcolor=\color{backcolour},   commentstyle=\color{codegreen},
  keywordstyle=\color{magenta},
  numberstyle=\tiny\color{codegray},
  stringstyle=\color{codepurple},
  basicstyle=\ttfamily\footnotesize,
  breakatwhitespace=false,         
  breaklines=true,                 
  captionpos=b,                    
  keepspaces=true,                 
  numbers=left,                    
  numbersep=5pt,                  
  showspaces=false,                
  showstringspaces=false,
  showtabs=false,                  
  tabsize=2
}
\lstset{style=mystyle}
\newcommand{\lfa}{\Leftrightarrow}
\newcommand{\N}{\mathbb{N}}

\title{Compilation de sujets de Colles de maths \textendash~Sup}
\author{ }
\date{ }

\usepackage{fancyhdr}
\pagestyle{fancy}
\renewcommand{\headrulewidth}{1pt}
\fancyhead{}
\fancyhead[L]{Sujets de colles}
\fancyhead[R]{\leftmark}
\newcommand{\R}{\mathbb{R}}
\newcommand{\abs}[1]{\lvert#1 \rvert}
\begin{document}
\maketitle
Un petit florilège des exercices de colle qui sont tombés en khôlle de mathématiques en première année de CPGE MPSI, il y à un peu de tout et de différents niveaux (qui ne sont pas encore indiqués pour l'instant, peut être un jour...)

Parfois, la solution des exercices n'est pas mise, en général parce que l'exercice était assez simple (ou que j'ai la flemme)


\tableofcontents
\section{Applications et fonctions}
\subsection{Bijections, injections etc etc}

\begin{exercise}[Une bijection de $\mathbb{R} \to { ]-1, 1[ }$\\]
  Soit $f:\R \to \R$, $f(x) = \frac{x}{1+\abs{x}}$
  \begin{enumerate}

    \item Montrer que $f$ est bijective.
  \item Montrer que $f(\R) \subset ]-1,1[$.
  \item Soit $y \in]-1,1[$. Résoudre $f(x) = y$. En déduire que l'application $g:\R \mapsto ]-1,1[$ induite par $f$ est bijective et expliciter $g^{-1}$.
\item Déterminer $f^{-1}\left(\left[0, \frac{3}{2}\right]\right)$ et $f\left(\rbrack\infty,0\rbrack\right)$

  \end{enumerate}
\end{exercise}

\begin{exercise}[Images directes et réciproques\\]
  Soit $f:E\to F$, $A,B\subset E$ et $C,D \subset F$
  \begin{enumerate}
    \item Montrer que $f(A\cap B) \subset f(A)\cap f(B)$ et que $f^{-1}(C\cap D) = f^{-1}(C)\cap f^{-1}(D)$.
    \item Donner un exemple où l'inclusion de la question 1 est stricte. Montrer qu'il y à égalité pour tous A,B si $f$ est injective.
    \item Réciproquement, on suppose que $f(A\cap B) = f(A)\cap f(B)$ pour tous $A,B\subset E$. Montrer que $f$ est injective (prendre $A=\{x\}$, $B=\{x'\}$)
    \item Montrer que $f(f^{-1}(C)) = C\cap f(E)$ et que $A\subset f^{-1}(f(A))$.
  \end{enumerate}
\end{exercise}
\begin{exercise}[Propriétés de relations binaires\\]
  On considère sur $\R$ les relation $x\mathcal{R}_1y \lfa \abs{x-y}\leq1$, $x\mathcal{R}_2y \lfa x-y \in \mathbb{Q}$, $x\mathcal{R}_3y \lfa \abs{x}\leq\abs{y}$.
  \begin{enumerate}

    \item Pour chacune, étudier réflectivité, symétrie, antisymétrique et transitivité.
    \item Pour $\mathcal{R}_2$, décrire la classe de $0$ et celle de $\sqrt{2}$ (on admet que $\sqrt{2} \in \mathbb{Q}$). Ces classes sont-elles égales?
    \item Pour $\mathcal{R}_3$, définir $x \sin y \lfa x\mathcal{R}_3y \land y \mathcal{R}_3 x$. Montrer que $\sin$ est une relation d'équivalence et décrire ses classes.
  \end{enumerate}

\end{exercise}

\subsection{Fonctions usuelles}

\begin{exercise}
  Soient $a, b \in \R^*_+, a\leq b$
  Montrer que:
  \[\ln \left( 1 + \frac{a}{b}\right)\ln\left(1+\frac{b}{a}\right) \leq \ln {(2)}^2\]

\end{exercise}
\begin{solution}
  Il faut poser la fonction $f:x\mapsto \frac{\ln(1+ax)}{\ln(1+bx)}$, puis étudier ses variations (elle est croissante) puis utiliser la croissance pour avoir $f\left(\frac 1 a\right)\geq f\left(\frac 1 b\right)$ puis déduire ce qu'il faut pour tuer l'exercice.

\end{solution}

\subsection{Equations fonctionnelles}
\begin{exercise}
  Trouver les $f\in\R^{\R^*}$ qui vérifient l'équation fonctionnelle:
  \[f(x)+3f\left(\frac 1 x\right) = x^2\]
\end{exercise}
\begin{solution}
  Raisonnement par analyse synthèse.\\
  Remplacer $x$ par $\frac 1 x$ puis substituer ce dont on a besoin et conclure (bien penser à faire la synthèse)

\end{solution}
\section{Somme et produits}
\begin{exercise}

  Calculer:
  \[\sum_{\substack{k=0\\k\equiv0 \ [3]}}^n \binom{n}{k}\]
\end{exercise}

\begin{solution}
  \[S_1 = {(1+1)}^n = 2^n = \sum_{k=0}^n\binom{n}{k}\]
  \[S_j = {(1+j)}^n = e^{\frac{ni\pi}{3}} = \sum_{k=0}^n\binom{n}{k} j^k\]
  \[S_{\bar{\jmath}} = {(1+\bar{\jmath})}^n = e^{-\frac{ni\pi}{3}} = \sum_{k=0}^n\binom{n}{k} \bar{\jmath}^k\]

  \[S_1+S_j+S_{\bar{\jmath}} = \sum_{k=0}^n(1+j+\bar{\jmath})\binom{n}{k}\]

  On remarque que si $k\equiv 0 \ [2]$, $ (1+j+\bar{\jmath}) = 3$, sinon $ (1+j+\bar{\jmath}) = 0$
  on a donc:

  \[S_1+S_j+S_{\bar{\jmath}} = 3\sum_{\substack{k=0\\k\equiv0 \ [3]}}^n \binom{n}{k}\]

  Et donc:
  \[\sum_{\substack{k=0\\k\equiv0 \ [3]}}^n \binom{n}{k} = \frac{2^n+e^{\frac{ni\pi}{3}}+e^{-\frac{ni\pi}{3}}}{3}\]

\end{solution}

\begin{exercise}

  Calculer la somme suivante: \[\sum_{1\leq i,j\leq n} \lvert i - j \rvert\]
\end{exercise}

% fin somme
\section{Arithmétique et coefficients binomiaux}
Il n'y aura surement qu'un exercice ici, il est assez inclassable. 
\begin{exercise}

  Montrer que $n$ est premier si et seulement si $\forall k \in \llbracket 1~;~n-1\rrbracket$, $n\mid\binom{n}{k}$
\end{exercise}

\begin{solution}

  Sens $\Rightarrow$ plutôt facile, bien penser a dire que n est premier pour pouvoir le factoriser

  Sens $\Leftarrow$:
  Faire par l'absurde en supposant que $n$ est composé puis voir que $n\nmid\binom{n}{k}$

\end{solution}

\section{Trigonométrie}
\subsection{Calcul de valeur trigonométrique }%
\label{sub:Calcul de valeur trigonométrique }

\begin{exercise}

  Calculer:
  \[\cos\left(\frac{\pi}{7}\right)\cos\left(\frac{2\pi}{7}\right)\cos\left(\frac{3\pi}{7}\right)\]
\end{exercise}

\begin{solution}

  Multiplier par $\sin\left(\frac{\pi}{7}\right)$ et remarquer que $\sin\left(\frac{3\pi}{7}\right) = \sin\left(\frac{4\pi}{7}\right)$

\end{solution}

\begin{exercise}
    calculer $\cos \left( \frac{\pi}{10} \right) $
\end{exercise}
\begin{solution}
    D'abord, il faut dé linéariser $\cos 5x$ puis évaluer l'expression obtenue en $\frac{\pi}{10}$ et l'on obtient un polynôme de degré 4 que l'on peut réécrire en un polynôme de degré 2 que l'on peut aisément résoudre.
\end{solution}
\subsection{Avec des fonctions trigonométrique à étudier}%
\label{sub:Avec des fonctions trigonométrique à étudier}

\begin{exercise}

  Soit $f:x \mapsto \arcsin\left(2x\sqrt{(1-x^2)}\right)$
  \begin{itemize}
    \item Donner son ensemble de définition et de dérivation
    \item Trouver $f'(x)$
    \item En déduire une simplification de $f$ en particulier la constante de différence avec une fonction que l'on reconnaitra.
  \end{itemize}
\end{exercise}

\begin{solution}
  \leavevmode
  \begin{itemize}
    \item Ensemble de définition : assez trivial, ($\arcsin$ défini sur $[-1, 1]$)
  
    \item  Ensemble de dérivation, chercher les moments où $ 2x\sqrt{1-x^2}=1$ (il y à trois intervalles différents)

    \item  Dérivée, choisir un des intervalles où la fonction est dérivable et dérivée avec composée de fonction, on repèrera une fonction dont la dérivée est la même. Deux fonctions dont les dérivées sont égales sont égales modulo une constante il faut donc chercher cette constante dans un certain intervalle où $f$ est dérivable.
\end{itemize}
\end{solution}


\section{Nombres complexes}%
\label{sec:Nombres complexes}

\begin{exercise}
    Déterminer $\sup_{\lvert z \lvert \leq 1} \lvert z^3+2iz\lvert$
\end{exercise}
\begin{solution}
    \begin{itemize}
        \item $\forall z \in \mathbb{C}$ tel que $\lvert z \lvert \leq 1$:
    \begin{align*}
        \lvert z^3+2iz \lvert &\leq \lvert z^3 \lvert +\lvert 2iz \lvert \\
                              &\leq \lvert z^3 \lvert +2 \lvert z \lvert 
    \end{align*}
    donc $\lvert z^3+2iz \lvert \leq 3$ par inégalité triangulaire
\item On cherche si on peut obtenir une égalité au lieu d'une inégalité, on peut intuiter assez aisément le maximum, et on a bien que pour ce nombre $\lvert z^3+2iz \lvert =3$
    \end{itemize}
\end{solution}

\section{Suites}%
\label{sec:Suites}
\subsection{Limites de suites}%
\begin{exercise}
    Étudier la suite définie par: \[
        u_0 \in \mathbb{R}_+ \land \forall n \in~\mathbb{N} u_{n+1} = 1+\frac{1}{4}u_n^2
    \]
\end{exercise}

\begin{exercise}
    Soit $(u_n=$ une suite décroissante de limite nulle. On pose: \[
        S_n = \sum_{k=1}^n (-1)^k u_k
    \]
    Étudier le comportement de la suite $(S_n)$.
\end{exercise}
\begin{solution}
    On pourra s'intéresser à $(S_{2n})$ et $(S_{2n+1})$. (On verra que ces deux suites sont adjacentes et après on est heureux :D).
\end{solution}
\label{sub:Limites de suites}
\begin{exercise}
    Soit $(n_n)_{n\in \mathbb{N}}$ une suite réelle bornée vérifiant: \[
        \forall n \in~\mathbb{N}^* a_{n+1} \leq \frac{1}{3}a_n + \frac{2}{3} a_{n-1} 
    \]

    \begin{enumerate}
        \item Montrer que $\left( a_{n+1} + \frac{2}{3}a_n \right) _{n \in~\mathbb{N}}$ est convergente. On note $\lambda$ sa limite.
        \item Soit $\epsilon \ge 0$. Pour $n \in~\mathbb{N}$, on pose $u_n = a_n - \frac{3}{5}\lambda$. Montrer qu'il existe $n_0 \in \mathbb{N}$ tel que l'on ait $\lvert u_{n+1} \lvert \frac{2}{3}\lvert u_n \lvert + \epsilon$ pour tout $n\geq n_0$.
        \item Étudier la convergence de la suite $(a_n)_{n\in\N}$.
    \end{enumerate}
\end{exercise}

\begin{exercise}
    D'après la formule de Viete: \[
        \sqrt{2+\sqrt{2+\sqrt{2+{\sqrt{2+\sqrt{\cdots}}}}}} = 2
    \]
    \begin{enumerate}
        \item Réécrire cette formule et montrer avec la notion de suite.
        \item Pour $\alpha \in \mathbb{R}^*_+$, que vaut $\sqrt{\alpha+\sqrt{\alpha+\sqrt{\alpha+{\sqrt{\alpha+\sqrt{\cdots}}}}}}$ ?
    \end{enumerate}
\end{exercise}

\begin{exercise}
    On pose pour tout entier naturel $n$ : $u_n = \sqrt{n} - \lfloor \sqrt{n} \rfloor$
    \begin{enumerate}
        \item Étudier $\lim u_{n^2+n}$ et en déduire que $(u_n)$ ne converge pas.
    \item Soient $a$ et $b$ deux entiers naturels $b\neq 0 \land a \le b$. Démontrer que $lim_{n^2b^2+2an} = \frac{a}{b}$
    \item Démontrer que tout élément de $\left[ O;1 \right]$ est limite d'une suite extraite de $(u_n)$.
    \end{enumerate}
\end{exercise}

\section{Ensembles}%
\label{sec:Ensembles}

\begin{exercise}
    
    \begin{enumerate}
        \item Soit $A$ une partie non vide et bornée de $\mathbb{R}$. On note $B=\left\{ \lvert x-y \lvert ,x,y \in A \right\} $.
        \item On pose $\delta(A) = \sup B$. Prouver que $\delta(A) = \sup A - \sup B$ 
    \end{enumerate}
\end{exercise}


\end{document}
